\section{讨论}
\label{sec:discussion}

\subsection{伦理与责任}
\label{sec:ethics}

系统把最终决定留给人工审阅者，这是有意的设计选择。自动拒绝会节省时间，但也会把
错误的责任藏进流程；在人机协同的审阅里，责任必须能落到一个具体的人和一次具体的
决定上。审计记录因此不是附加功能，而是系统的一部分：它记录模型提议、审阅者动作、
时间与理由，使事后复核成为可能。

我们遵循数据最小化原则：评估只使用合成文档；如果部署到真实场景，应先确认数据授权、
去标识化方式与保留期限。模型文档与数据说明应随系统一起维护\cite{mitchell2019,gebru2021}，
而不是等审计时才补。

\subsection{局限}
\label{sec:limitations}

评估有三个明显局限。第一，语料库是合成的、规模很小，主题只有两个，时间节省不一定
能推广到真实文档。第二，只有两名审阅者，审阅者间的差异没有单独分析。第三，评估没有
覆盖长期使用中的自动化偏差：审阅者可能逐渐减少独立判断。下一步需要更大、更多样的
语料库，以及一段时间的纵向观察。
